\documentclass[11pt,a4paper]{article}
\usepackage[T1]{fontenc}
\usepackage[utf8]{inputenc}
\usepackage{lmodern,microtype}
\usepackage[a4paper,margin=26mm,headheight=15pt]{geometry}
\usepackage{amsmath,amssymb,amsthm,mathtools,bm,mathrsfs}
\usepackage{booktabs,array,longtable,enumitem}
\usepackage{xcolor,xurl,fancyhdr,hyperref}
\hypersetup{colorlinks=true,linkcolor=blue!40!black,citecolor=blue!40!black,
 urlcolor=blue!40!black,pdftitle={Algebraic Convergence Loci for Nonlinear Hahn--Fuchsian Systems},
 pdfauthor={Research article prepared for Vladimir Reshetnikov}}
\urlstyle{same}
\numberwithin{equation}{section}
\newtheorem{theorem}{Theorem}[section]
\newtheorem{proposition}[theorem]{Proposition}
\newtheorem{lemma}[theorem]{Lemma}
\newtheorem{corollary}[theorem]{Corollary}
\theoremstyle{definition}
\newtheorem{definition}[theorem]{Definition}
\newtheorem{example}[theorem]{Example}
\theoremstyle{remark}
\newtheorem{remark}[theorem]{Remark}
\newcommand{\C}{\mathbb C}
\newcommand{\R}{\mathbb R}
\newcommand{\N}{\mathbb N}
\newcommand{\G}{\Gamma}
\newcommand{\HH}{\mathscr H}
\newcommand{\mm}{\mathfrak m}
\newcommand{\Res}{\mathcal R}
\newcommand{\Anc}{\mathcal A}
\newcommand{\cF}{\mathcal F}
\newcommand{\cV}{\mathcal V}
\newcommand{\cW}{\mathcal W}
\newcommand{\supp}{\operatorname{supp}}
\newcommand{\val}{\operatorname{v}}
\newcommand{\spec}{\operatorname{spec}}
\newcommand{\rank}{\operatorname{rank}}
\newcommand{\diag}{\operatorname{diag}}
\newcommand{\wt}{\operatorname{wt}}
\newcommand{\norm}[1]{\left\lVert #1\right\rVert}
\newcommand{\abs}[1]{\left\lvert #1\right\rvert}
\newcommand{\snap}{\texttt{3d5973524506411392a911470b5ddc35521568ea}}
\newcommand{\proj}{\operatorname{pr}}
\setlength{\parskip}{3pt}
\setlength{\parindent}{1.15em}
\setlist{itemsep=3pt,topsep=5pt,leftmargin=1.7em}
\linespread{1.025}
\pagestyle{fancy}
\fancyhf{}
\fancyhead[L]{\small Nonlinear Hahn--Fuchsian systems}
\fancyhead[R]{\small Algebraic convergence loci}
\fancyfoot[C]{\thepage}
\renewcommand{\headrulewidth}{0.3pt}
\setcounter{tocdepth}{1}
\begin{document}
\hypersetup{pageanchor=false}
\begin{titlepage}
\thispagestyle{empty}
\vspace*{8mm}
\noindent{\large\scshape Transseries research}\par
\vspace{10mm}
\noindent{\LARGE\bfseries Algebraic Convergence Loci\par}
\vspace{3mm}
\noindent{\LARGE\bfseries for Nonlinear\par}
\vspace{3mm}
\noindent{\LARGE\bfseries Hahn--Fuchsian Systems\par}
\vspace{7mm}
\noindent{\large Finite resonances, bounded-exponent obstructions,\par
and universality with accumulating supports}
\vspace{9mm}\hrule\vspace{6mm}
\noindent Research article prepared for Vladimir Reshetnikov\par
\noindent 29 September 2026
\vspace{7mm}
\begin{abstract}
Consider $Dy=Ay+F(x,y)$, where $D=x\,d/dx$, $A$ is an arbitrary complex
matrix, $F(0,0)=0$, $F_y(0,0)=0$, and the coefficients of $F$ are absolutely
convergent Hahn series supported in a fixed well-ordered positive real
exponent monoid. Finite accumulation of exponents is allowed.

We prove that the positive, logarithm-free formal solutions are parametrized
by a finite affine algebraic set of resonant coefficients. More importantly,
\emph{the subset of parameters producing absolutely convergent solutions is
itself affine algebraic}, even when the spectral separation is zero.
The proof separates a bounded exponent window from a contractive analytic
tail, then takes a finite-dimensional algebraic quotient of coefficient
families by $\ell^1$. The additional convergence equations have strictly
bounded weighted degree. Convergent solutions have a common positive
radius on each compact subset of their algebraic parameter locus.

This conclusion is sharp in a geometric sense: every affine algebraic set
occurs, after a harmless normalization, as the convergence locus of a
triangular system with \emph{no formal compatibility obstruction}.
We also prove an exact universal spectral-gap criterion, a finite-ancestry
bound for formal obstructions, and a sharp nonlinear example in which the
forcing stays away from the spectrum but quadratic interactions create
an accumulation of small divisors. Its convergence threshold is exactly
$p>2$.

All theorem-level arguments are given in the article. The accompanying
program checks 670 exact scalar equalities, including 21 finite systems.
These computations are consistency checks, not a formal verification of
the infinite-support results.
\end{abstract}
\vfill
\noindent\small\textbf{Status.} A self-contained research extension of selected
ProveIt material. Historical priority has not been independently established;
no claim of a solution to a famous general transseries conjecture is made.\par
\vspace{2mm}
\noindent\small\textbf{Inspected repository snapshot.}\quad\snap
\end{titlepage}
\hypersetup{pageanchor=true}
\tableofcontents
\newpage

\section{The question and the proposed advance}
\label{sec:intro}

Formal solvability, the absence of logarithms, and analytic convergence are
three different questions. For ordinary regular-singular equations their
relationship is familiar. With Hahn exponents, infinitely many terms can
lie below one fixed spectral value, and formal division can destroy
absolute convergence without changing the support.

The ProveIt package \emph{Finite Resonance Certificates and Sharp Analytic
Normalization for Hahn--Fuchsian Systems} treats a linear matrix-gauge
problem. It identifies a universal separation condition involving spectral
differences, and gives accumulation examples where a formal normalizer
diverges \cite{repo-linear}. Its further directions ask about constrained
supports, cancellation, and the geometry of exceptional parameter sets.
Those results are not reproved here as a new discovery. We change the
problem: the unknown is a nonlinear solution germ, not a linear gauge, and
we classify \emph{which resonant parameters yield convergent germs}.

The precise research question is:
\begin{quote}
Can infinitely many small-divisor conditions impose a genuinely
nonalgebraic convergence condition on the finite resonant parameters of a
positive Hahn solution of a nonlinear regular-singular equation?
\end{quote}
Within the class specified in Section~\ref{sec:framework}, the answer is
\emph{no}. The convergence locus is algebraic. Conversely, it can be any
affine algebraic set, so there is no general smoothness, linearity, or
irreducibility theorem to replace this answer.

\subsection{A guide to the results}
There are two finite-dimensional obstructions, with very different meanings.
The formal compatibility map $\Psi$ depends on finitely many input
coefficients. The analytic map $\chi$ has finitely many polynomial
components, but can depend on the summability of infinite coefficient
families. The main description is
\begin{equation}\label{eq:main-description}
 \cV_{\mathrm{formal}}=\{c:\Psi(c)=0\},\qquad
 \cV_{\mathrm{an}}=\{c:\Psi(c)=0,\ \chi(c)=0\}.
\end{equation}
The second equality is not a finite-coefficient decision algorithm.

Theorem~\ref{thm:formal} constructs $\Psi$ without an analytic assumption.
Theorem~\ref{thm:window} locates every possible convergence obstruction in
one bounded exponent window. Theorem~\ref{thm:algebraic} converts that
infinite window into the finite polynomial map $\chi$.
Theorem~\ref{thm:universality} realizes arbitrary affine algebraic sets as
$\cV_{\mathrm{an}}$ while $\cV_{\mathrm{formal}}$ is affine space.
Theorem~\ref{thm:gap} characterizes when $\chi$ vanishes for every input.
The example in Section~\ref{sec:hidden} shows why checking only the original
forcing support misses nonlinear small divisors.

\subsection{Relation to established theory}
Generalized-series fixed-point methods are established. Van der Hoeven's
operator theory includes an implicit-function theorem for Noetherian
operators \cite{vdh}. We use a direct positive-valuation argument in the
restricted setting needed here. Neumann support finiteness is also
classical; the inspected repository already has a Lean development of
finite ordered-word representations \cite{repo-words}.

Convergence of generalized series satisfying analytic differential equations
has substantial antecedents. Gontsov and Goryuchkina study complex-exponent
series with real parts tending to infinity, and prove a majorant convergence
theorem \cite{gg2015}; their later work develops a broader functional-equation
framework \cite{gg2024}. Our accumulating coefficient supports are not an
ordinary holomorphic power series in $x$ at zero and are not covered merely
by identifying our equation as regular singular. Absolute convergence for
Hahn series itself is a standard notion, also used in higher-rank tropical
geometry \cite{js2023}.

The proposed advance is the algebraicity and universality of the
\emph{parameter-dependent absolute-convergence locus at zero spectral gap}.
The proof is supplied independently below. A targeted literature and
repository comparison did not establish historical priority for this exact
statement; the article does not assert exhaustive coverage of either.

\subsection{Scope}
Exponents are real and nonnegative, the exponent monoid is well ordered,
and solutions have positive valuation. The constant matrix need not be
semisimple and need not have real spectrum. We consider no powers of
$\log x$ in the solution ansatz. Failure of formal compatibility means
failure in this logarithm-free class, not nonexistence of other solutions.
Divergence means failure of absolute Hahn convergence, not nonexistence of
an analytic function with some weaker asymptotic interpretation.

Under $x=e^{-t}$, the monomials become $e^{-\gamma t}$ and $D=-d/dt$.
Thus this is also a theorem about exponential transseries sectors with
possibly accumulating positive actions. It is not a theorem about every
logarithmic or exponential level of the full transseries field.

\section{The coefficient algebra and the nonlinear model}
\label{sec:framework}

\subsection{A fixed well-ordered monoid}
Let $\G\subset\R_{\geq0}$ be a nontrivial additive submonoid, containing
$0$, and well ordered in the usual real order. Put
\begin{equation}\label{eq:delta}
 \delta=\min\G_{>0}>0,
 \qquad
 \HH_\G=\left\{\sum_{\gamma\in\G}u_\gamma x^\gamma:
 u_\gamma\in\C\right\},
 \qquad
 \mm_\G=\{u:u_0=0\}.
\end{equation}
Every subset of $\G$ is well ordered. We do not require
$\G\cap[0,L]$ to be finite.

\begin{lemma}[Finite additive ancestry]\label{lem:ancestry}
For every $\rho\in\G$, the set
\begin{equation}\label{eq:ancestors}
 \Anc(\rho)=\{\alpha\in\G:\rho-\alpha\in\G\}
\end{equation}
is finite. If $\beta\in\Anc(\rho)$ and
$\alpha\in\Anc(\beta)$, then $\alpha\in\Anc(\rho)$.
Every ordered decomposition of $\rho$ into positive elements of $\G$
has length at most $\lfloor\rho/\delta\rfloor$, and there are only
finitely many such decompositions.
\end{lemma}
\begin{proof}
If $\Anc(\rho)$ were infinite, successively taking least unused elements
would give a strictly increasing sequence $\alpha_n$ in that set. Then
$\rho-\alpha_n$ would be a strictly decreasing sequence in $\G$, contrary
to well ordering. The hereditary assertion follows from
$\rho-\alpha=(\rho-\beta)+(\beta-\alpha)\in\G$.
Every positive summand is at least $\delta$, which proves the length bound.
For each fixed length, finiteness follows inductively from the two-summand
assertion just proved. Only finitely many lengths are possible.
\end{proof}

The same two-set argument proves that the sum of two well-ordered subsets
of $\R_{\geq0}$ is well ordered and has finite coefficient antidiagonals.
Consequently, a well-ordered set $E\subset\R_{>0}$ generates a well-ordered
monoid: below a fixed real bound only finitely many word lengths occur,
and each fixed-length sumset is well ordered. This supplies many admissible
examples. The stronger general ordered-group word lemma in
\cite{repo-words} is not claimed as a new result here.

Convolution therefore defines a ring on $\HH_\G$. Define $D$ strongly by
\begin{equation}\label{eq:D}
 D\!\left(\sum_\gamma u_\gamma x^\gamma\right)
   =\sum_\gamma\gamma u_\gamma x^\gamma.
\end{equation}
Finite convolutions give the Leibniz rule. For vector series use
$\val(y)=\min\supp(y)$, with $\val(0)=+\infty$.

\subsection{The equation}
Fix $A\in\operatorname{Mat}_m(\C)$. Consider
\begin{equation}\label{eq:system}
 Dy=Ay+F(x,y),\qquad y\in\mm_\G^m,
\end{equation}
where
\begin{equation}\label{eq:F}
 F(x,y)=\sum_{\alpha\in\G}\sum_{k\in\N^m}
                F_{\alpha,k}x^\alpha y^k,
 \qquad F_{\alpha,k}\in\C^m,
 \qquad F_{0,0}=0,\quad F_{0,e_i}=0.
\end{equation}
Here $y^k=y_1^{k_1}\cdots y_m^{k_m}$ and $|k|=\sum k_i$.
The last two conditions are exactly the formal interpretation of
$F(0,0)=0$ and $F_y(0,0)=0$. Constant and linear terms with \emph{positive}
$x$-exponent are allowed.

For fixed $\gamma$, a term contributing to $[x^\gamma]F(x,y)$ has
$|k|\leq\gamma/\delta$. Every $x$-exponent and every factor exponent in
that term belongs to $\Anc(\gamma)$. Thus the coefficient is a finite
sum. If $|k|=1$, then $\alpha>0$; if $\alpha=0$, then $|k|\geq2$.
In either case every solution coefficient used on the right has exponent
strictly less than $\gamma$.

\begin{lemma}[Uniform formal valuation gain]\label{lem:gain}
For $y,z\in\mm_\G^m$,
\begin{equation}\label{eq:gain}
 \val\bigl(F(x,y)-F(x,z)\bigr)\geq\val(y-z)+\delta.
\end{equation}
\end{lemma}
\begin{proof}
Factor the difference of each monomial $y^k-z^k$ into terms with one
factor from $y-z$. A positive $x$-exponent supplies a gain of at least
$\delta$. Otherwise $|k|\geq2$, and at least one additional positive
solution factor supplies that gain. Coefficientwise sums are finite, so
cancellation can only increase valuation.
\end{proof}

\subsection{Absolute Hahn convergence}
Choose any norm on $\C^m$; for estimates below take the maximum-coordinate
norm and its induced matrix norm. Define
\begin{equation}\label{eq:norm}
 \norm{y}_r=\sum_{\gamma\in\G}\norm{y_\gamma}r^\gamma,
 \qquad B_r=\ell^1(\G,r^\gamma).
\end{equation}
Scalar $B_r$ is a Banach convolution algebra. Write $B_r^+$ for its
positive-support ideal. A formal series is called \emph{convergent} if it
belongs to some $B_r$, $r>0$.

The analytic assumption on $F$, imposed from Section~\ref{sec:window}
onward, is that for some $r_0,s_0>0$,
\begin{equation}\label{eq:analytic-F}
 \sum_{\alpha,k}\norm{F_{\alpha,k}}r_0^\alpha s_0^{|k|}<\infty.
\end{equation}
This is a precise joint majorant condition; it is not inferred from formal
support. Put
\begin{align}
 b(r)&=\sum_\alpha\norm{F_{\alpha,0}}r^\alpha,\label{eq:b}\\
 L(r,s)&=\sum_{\alpha,\,|k|\geq1}
   |k|\norm{F_{\alpha,k}}r^\alpha s^{|k|-1}.
 \label{eq:L}
\end{align}
For $s<s_0$ and $r\leq r_0$ these are finite, and
\begin{equation}\label{eq:smallness}
 b(r)\longrightarrow0\quad(r\downarrow0),\qquad
 L(r,s)\longrightarrow0\quad(r,s\downarrow0).
\end{equation}
The assertions follow by dominated convergence, using the vanished
constant-linear coefficients and, for $L$, a fixed smaller $y$-radius to
absorb the factor $|k|$. For $\norm{y}_r,\norm{z}_r\leq s$,
\begin{equation}\label{eq:lipschitz}
 \norm{F(x,y)-F(x,z)}_r\leq L(r,s)\norm{y-z}_r,
 \qquad \norm{F(x,y)}_r\leq b(r)+L(r,s)\norm{y}_r.
\end{equation}
These elementary inequalities supply every analytic fixed-point estimate
used below.

\section{Finite algebraic parametrization of formal solutions}
\label{sec:formal}

\subsection{Kernel and cokernel splittings}
Let
\begin{equation}\label{eq:resonance}
 \Res=\G_{>0}\cap\spec(A),\qquad T_\gamma=\gamma I-A.
\end{equation}
Since $\gamma$ is real, only positive real eigenvalues lying in $\G$ occur.
For $\gamma\notin\Res$ put $R_\gamma=T_\gamma^{-1}$ and $Q_\gamma=0$.
For $\rho\in\Res$, choose complements
\[
 \C^m=\ker T_\rho\oplus U_\rho,
 \qquad \C^m=\operatorname{im}T_\rho\oplus W_\rho.
\]
Let $Q_\rho$ be projection onto $W_\rho$, and let $R_\rho$ be the inverse
of $T_\rho:U_\rho\to\operatorname{im}T_\rho$, extended by zero on
$W_\rho$. Then
\begin{equation}\label{eq:splitting}
 T_\rho R_\rho=I-Q_\rho,\qquad
 P_\rho:=I-R_\rho T_\rho
 \text{ projects onto }\ker T_\rho.
\end{equation}
A kernel and a cokernel need not be the same subspace. This distinction is
essential for Jordan blocks.

Choose bases of all $\ker T_\rho$. Let $c\in\C^q$ collect their
coordinates, where
\begin{equation}\label{eq:q}
 q=\sum_{\rho\in\Res}\dim\ker T_\rho\leq m.
\end{equation}
Denote the corresponding kernel vector at $\rho$ by $C_\rho(c)$, and put
$C(x;c)=\sum_{\rho\in\Res}C_\rho(c)x^\rho$. Give each parameter $c_i$
the positive weight $\tau_i=\rho$ of its resonance.

\subsection{The universal formal candidate}

\begin{theorem}[Formal compatibility is finite algebraic]
\label{thm:formal}
For every $c\in\C^q$ there is a unique formal candidate
$Y(c)\in\mm_\G^m$ satisfying
\begin{equation}\label{eq:candidate}
 Y=C(x;c)+\mathscr R F(x,Y),
 \qquad [x^\gamma]\mathscr R u=R_\gamma u_\gamma.
\end{equation}
Every coefficient $Y_\gamma(c)$ is a polynomial in $c$. There is a finite
polynomial map
\begin{equation}\label{eq:Psi}
 \Psi(c)=\bigl(Q_\rho[x^\rho]F(x,Y(c))\bigr)_{\rho\in\Res}
       \in\bigoplus_{\rho\in\Res}W_\rho
\end{equation}
such that $c\mapsto Y(c)$ is a bijection from $\Psi^{-1}(0)$ onto the
positive formal solutions of \eqref{eq:system}. Moreover:

\begin{enumerate}
\item Only the finite hereditary set
$\Anc_*:=\bigcup_{\rho\in\Res}\Anc(\rho)$ is needed to compute $\Psi$.
\item Only coefficients $F_{\alpha,k}$ with $\alpha\in\Anc_*$ and
$|k|\leq\lfloor R_*/\delta\rfloor$ can be needed, where
$R_*=\max\Res$ when $\Res\ne\varnothing$.
\item Each monomial $c^\nu$ occurring in $Y_\gamma$ has
$\wt(\nu):=\sum_i\tau_i\nu_i\leq\gamma$.
\end{enumerate}
For $\Res=\varnothing$ the map $\Psi$ is empty and there is one formal
solution.
\end{theorem}
\begin{proof}
Define $Y^{(0)}=0$ and
$Y^{(n+1)}=C+\mathscr R F(x,Y^{(n)})$.
The operator $\mathscr R$ cannot lower exponents. Lemma~\ref{lem:gain}
therefore gives
\[
 \val(Y^{(n+1)}-Y^{(n)})\geq(n+1)\delta.
\]
Every coefficient eventually stabilizes. Its limit is supported in $\G$,
hence is a Hahn series, and finite coefficientwise evaluation shows that
it satisfies \eqref{eq:candidate}. If two candidates have a nonzero
difference of least exponent $\beta$, \eqref{eq:gain} would make the same
difference have valuation at least $\beta+\delta$, a contradiction.

Equivalently, at each exponent one computes the finite polynomial
\begin{equation}\label{eq:recursion}
 f_\gamma(c)=[x^\gamma]F(x,Y_{<\gamma}(c)),\qquad
 Y_\gamma(c)=R_\gamma f_\gamma(c)+C_\gamma(c),
\end{equation}
where $C_\gamma=0$ off $\Res$. The right side uses only lower
coefficients. All dependencies lie in the finite set $\Anc(\gamma)$,
which proves polynomiality and finite dependence without making any
local-finiteness assumption on the whole low-exponent region.

Multiplying \eqref{eq:recursion} by $T_\gamma$ gives the exact residual
\begin{equation}\label{eq:residual}
 DY-AY-F(x,Y)=-\sum_{\rho\in\Res}\Psi_\rho(c)x^\rho.
\end{equation}
Thus $\Psi=0$ is sufficient. For any actual formal solution $y$, take
$C_\rho=P_\rho y_\rho$. Applying $R_\gamma$ to its coefficient equations
gives \eqref{eq:candidate}, so uniqueness proves necessity and the claimed
bijection.

Hereditary ancestry proves the first assertion. At exponents at most
$R_*$, a monomial has at most $R_*/\delta$ positive solution factors,
which proves the second. Finally, induction in \eqref{eq:recursion} shows
that the sum of the parameter weights in a product cannot exceed the sum
of its solution-factor exponents, which is at most the output exponent.
The free term at $\rho$ has exactly its assigned weight $\rho$.
\end{proof}

\begin{remark}[A finite certificate is not a bounded ordinary jet]
An infinite number of coefficients can lie below $R_*$. The theorem uses
only additive divisors of the \emph{exact} resonances. In particular,
``retain everything below the largest resonance'' is generally not a finite
algorithm. The finite set $\Anc_*$ can be much smaller.
\end{remark}

\begin{proposition}[Coordinate independence]\label{prop:intrinsic}
Different kernel/cokernel splittings give polynomially isomorphic formal
parameter varieties. Their subsets of convergent solutions correspond
under the same isomorphism.
\end{proposition}
\begin{proof}
On a formal solution, the new parameter at $\rho$ is the new kernel
projection applied to $Y_\rho(c)$. This is polynomial by
Theorem~\ref{thm:formal}. Reversing the splittings gives the inverse map,
since both parameter maps describe the same solution coefficients.
Convergence concerns the solution itself and is unchanged.
\end{proof}

\subsection{A sharp degree bound}
Treat the finitely relevant entries of $F_{\alpha,k}$ as variables as well
as the kernel coordinates $c$. The entries of $A,R_\gamma,Q_\gamma$ are
fixed constants.

\begin{proposition}[Total coefficient degree]\label{prop:degree}
The polynomial $Y_\gamma$ has total degree at most
$2\lfloor\gamma/\delta\rfloor-1$. The same bound holds for
$\Psi_\rho$ with $\gamma=\rho$. This bound is attained for arbitrarily
large integer $\rho/\delta$.
\end{proposition}
\begin{proof}
Expand the recursion into rooted substitution trees. A coefficient vertex
$F_{\alpha,k}$ has $|k|$ children and contributes exponent $\alpha$; a
free-parameter leaf contributes its resonance exponent. A zero-exponent
coefficient vertex has at least two children. Every other vertex has
positive exponent cost at least $\delta$. The sum of all costs is the
output exponent $\gamma$.

If $P$ is the number of positive-cost vertices and $Z$ the number of
zero-cost vertices, then $P\leq\lfloor\gamma/\delta\rfloor$.
Every leaf has positive cost, and a tree has at most its number of leaves
minus one vertices of arity at least two. Thus $Z\leq P-1$ and the number
of coefficient or parameter factors is at most $2P-1$.
Finite ancestry and the finite word-length bound make this a finite
polynomial expansion, not an unproved infinite tree sum.
The sharp family is established next.
\end{proof}

\begin{example}[Exact sharp obstruction]\label{ex:sharp}
For $n\geq2$, consider on $\G=\delta\N$
\begin{equation}\label{eq:sharp}
 (D-n\delta)y=a x^\delta+b y^2.
\end{equation}
The obstruction at $n\delta$ is exactly
\begin{equation}\label{eq:sharp-obstruction}
 \Psi_n=\frac{a^n b^{n-1}}{((n-1)!)^2\delta^{2n-2}}.
\end{equation}
Its total degree is $2n-1$, proving sharpness.

For completeness, set $z=x^\delta$, $\theta=z\,d/dz$, and, initially
for $b\ne0$, substitute $y=-(\delta/b)\theta u/u$, $u(0)=1$.
The Riccati equation becomes
\[
 \theta(\theta-n)u+qzu=0,\qquad q=ab/\delta^2.
\]
For $1\leq j<n$, its coefficients satisfy
$u_j=q u_{j-1}/(j(n-j))$, so
$u_{n-1}=q^{n-1}/((n-1)!)^2$.
The obstruction in the $y$ equation is $a u_{n-1}$, by multiplying the
linear residual by $-\delta^2/(bu)$ and comparing at the first obstructed
exponent. This is \eqref{eq:sharp-obstruction}.
Both sides are polynomial in $a,b$, so the identity extends to $b=0$.
\end{example}

\subsection{Choosing the monoid from the equation}
The fixed monoid is not necessarily an artificial restriction. Suppose
all coefficient exponents of $F$ lie in a given well-ordered positive
monoid $\G_0$. Adjoin all positive real eigenvalues of $A$ and take the
generated monoid $\G$. It is again well ordered. Every positive real Hahn
solution of the equation then lies in $\HH_\G^m$.

Indeed, read its finitely many kernel coefficients at the eigenvalues,
and construct the candidate in $\HH_\G^m$. Both satisfy the coefficientwise
fixed-point equation with those kernel data. Their two supports and the
input monoid have a common positive lower bound. Factoring their difference
as in Lemma~\ref{lem:gain} rules out a least nonzero coefficient of the
difference. Thus the arbitrary positive solution equals the constructed
candidate. This argument requires positive real support, not a general
higher-rank transseries support.

\section{All divergence is confined to a bounded exponent window}
\label{sec:window}

Assume \eqref{eq:analytic-F}. Define the spectral cutoff
\begin{equation}\label{eq:cutoff}
 H=\max\bigl(\{0\}\cup(\spec(A)\cap\R_{>0})\bigr).
\end{equation}
It includes positive real eigenvalues whether or not they belong to $\G$.
Write $J=\G\cap(0,H]$. This set may be infinite.

\begin{lemma}[The tail inverse is bounded]\label{lem:tail-inverse}
There is a finite constant
\begin{equation}\label{eq:tail-M}
 M_H=\sup_{\gamma\in\G,\ \gamma>H}\norm{(\gamma I-A)^{-1}}<\infty.
\end{equation}
\end{lemma}
\begin{proof}
The nonempty set $\G\cap(H,\infty)$ has a least element $\gamma_+>H$.
There is no real eigenvalue of $A$ in $[\gamma_+,\infty)$.
The resolvent is continuous and bounded on each compact subinterval of
that ray, and for $\gamma>\norm A$ the geometric-series estimate gives
$\norm{(\gamma I-A)^{-1}}\leq(\gamma-\norm A)^{-1}$.
Together these estimates prove the claim. The right-hand gap at $H$ follows
from well ordering; no left-hand gap has been assumed.
\end{proof}

For an explicitly certified norm bound one may instead use a larger cutoff
$H'>\norm A$, giving $M_{H'}\leq(H'-\norm A)^{-1}$.
The smaller cutoff \eqref{eq:cutoff} is useful for sharp degree bounds.

\begin{theorem}[Bounded-window convergence criterion]
\label{thm:window}
For every formal candidate $Y(c)$ of Theorem~\ref{thm:formal}, whether or
not $\Psi(c)=0$,
\begin{equation}\label{eq:window-equiv}
 Y(c)\in B_r\text{ for some }r>0
 \quad\Longleftrightarrow\quad
 \sum_{\gamma\in J}\norm{Y_\gamma(c)}<\infty.
\end{equation}
In particular, a compatible formal solution cannot acquire a new absolute
convergence obstruction above the largest positive real eigenvalue.
\end{theorem}
\begin{proof}
On a bounded interval of exponents, the positive weights $r^\gamma$ are
bounded above and below by positive constants. Thus membership of the
truncation in any $B_r$ is equivalent to the right side of
\eqref{eq:window-equiv}. Necessity follows at once.

For sufficiency put
$p=\sum_{\gamma\in J}Y_\gamma(c)x^\gamma$ and write $Y=p+z$.
The assumed coefficient sum is finite. For $r\leq1$,
\[
 \norm p_r\leq r^\delta\sum_{\gamma\in J}\norm{Y_\gamma(c)},
\]
so $\norm p_r\to0$. Let $\Pi_{>H}$ denote exponent truncation, and let
$\mathscr R_{>H}$ apply $(\gamma I-A)^{-1}$ on that tail.
The required tail fixed-point equation is
\begin{equation}\label{eq:tail-fixed}
 z=\mathscr R_{>H}\Pi_{>H}F(x,p+z).
\end{equation}
Both truncation and the inverse are bounded on $B_r$, with norms at most
$1$ and $M_H$.

Choose $0<\eta<s_0$. By first making $\eta$ small and then $r$ small,
we can arrange
\begin{equation}\label{eq:tail-conditions}
 \norm p_r\leq\eta/4,\qquad
 M_H b(r)\leq\eta/8,\qquad M_H L(r,\eta)\leq1/2.
\end{equation}
The map on the right of \eqref{eq:tail-fixed} sends the closed tail ball
$\norm z_r\leq\eta/2$ into itself: at $z=0$ its norm is at most
$\eta/8+\eta/8=\eta/4$, and its Lipschitz constant is at most $1/2$.
Banach's contraction theorem gives a convergent tail.

Its formal coefficients equal those of the original candidate. Above $H$
there is no resonance and both tails satisfy \eqref{eq:tail-fixed}.
If they differed, their least different exponent would gain at least
$\delta$ under the right side by Lemma~\ref{lem:gain}, which is impossible.
The finite possible formal residuals below $H$ do not enter this argument.
\end{proof}

\begin{corollary}[Locally finite exponent supports]\label{cor:locfinite}
If $\G\cap[0,L]$ is finite for every $L$, every positive formal solution
of \eqref{eq:system} is convergent. In particular this holds for monoids
generated by finitely many positive real numbers.
\end{corollary}
\begin{proof}
The set $J$ is then finite. A finitely generated positive real monoid is
locally finite because each generator count is bounded under a fixed
exponent cutoff.
\end{proof}

\subsection{Analytic evaluation}
A convergent Hahn series defines a holomorphic function on the logarithmic
cover of $0<|x|<r$. For $0<r'<r$,
\begin{equation}\label{eq:derivative-bound}
 \norm{Dy}_{r'}\leq
 \frac{\norm y_r}{e\log(r/r')},
\end{equation}
since $\sup_{t\geq0}t(r'/r)^t=1/(e\log(r/r'))$.
Thus termwise differentiation is valid on smaller radii. Absolute
convolution and \eqref{eq:analytic-F} also justify substitution into $F$.
A compatible convergent candidate consequently solves the actual
differential equation on that cover. These statements do not require
$F$ or $y$ to extend holomorphically across $x=0$ in the ordinary sense.

\section{The algebraic convergence-locus theorem}
\label{sec:algebraic}

This is the central result. The bounded exponent window can contain
infinitely many non-summable coefficients. Nevertheless its dependence
on finitely many positive-weight resonant parameters has bounded degree.
That finite polynomial complexity is sufficient.

\begin{theorem}[Algebraicity at zero spectral gap]
\label{thm:algebraic}
Assume \eqref{eq:analytic-F}. There are finitely many polynomials
$\chi_1,\ldots,\chi_s\in\C[c_1,\ldots,c_q]$ such that
\begin{equation}\label{eq:analytic-variety}
 \{c:Y(c)\text{ is convergent}\}=\{c:\chi(c)=0\}.
\end{equation}
Consequently the positive convergent solutions are parametrized by the
affine algebraic set $\cV_{\mathrm{an}}=V(\Psi,\chi)$.
If $H>0$, the polynomials can be chosen to use only monomials satisfying
\begin{equation}\label{eq:strict-weight}
 \sum_{i=1}^q\tau_i\nu_i<H.
\end{equation}
The number $s$ is at most the number of such monomials. If $q>0$, their
ordinary degrees are at most $\lceil H/\min_i\tau_i\rceil-1$.
If $H=0$, there is no analytic obstruction.
\end{theorem}
\begin{proof}
Suppose first that $H>0$. Theorem~\ref{thm:formal} gives, on the entire
bounded window $J$, a uniform finite monomial expansion
\begin{equation}\label{eq:finite-monomials}
 Y_\gamma(c)=\sum_{\nu:\,\wt(\nu)\leq H} a_{\gamma,\nu}c^\nu,
 \qquad a_\nu=(a_{\gamma,\nu})_{\gamma\in J}.
\end{equation}
There are only finitely many allowed $\nu$, since all parameter weights
are positive. For $q=0$ the one possible monomial is the constant $1$.

Let $\mathcal E=(\C^m)^J$ be the vector space of all coefficient families
on $J$, and let $\mathcal L=\ell^1(J;\C^m)$ be its linear subspace of
absolutely summable families. Form the \emph{algebraic} quotient
$\mathcal E/\mathcal L$; no norm or Hausdorff topology on this quotient
is used. The classes $[a_\nu]$ span a finite-dimensional subspace $W$.
Choose a basis $w_1,\ldots,w_s$ of $W$ and write
\[
 [a_\nu]=\sum_{j=1}^s q_{j\nu}w_j.
\]
Define
\begin{equation}\label{eq:chi}
 \chi_j(c)=\sum_{\nu}q_{j\nu}c^\nu.
\end{equation}
The family $(Y_\gamma(c))_{\gamma\in J}$ lies in $\mathcal L$ exactly
when its quotient class vanishes, and by linear independence this is
exactly $\chi_1(c)=\cdots=\chi_s(c)=0$.
Theorem~\ref{thm:window} proves \eqref{eq:analytic-variety}.

If $\wt(\nu)=H$, then $a_{\gamma,\nu}=0$ for $\gamma<H$ by weighted
degree. Its only possible nonzero entry in $J$ is at $\gamma=H$.
Such a family is summable, so $[a_\nu]=0$. Those monomials may therefore
be omitted, giving the strict inequality \eqref{eq:strict-weight}.
The equation-count and ordinary-degree bounds follow immediately.
When $H=0$, the set $J$ is empty, and Theorem~\ref{thm:window} makes every
candidate convergent.
\end{proof}

\begin{remark}[Why this is more than a spectral-gap statement]
No lower bound on the nonzero denominators below $H$ was used. Infinite
families of unbounded coefficients may cancel after specializing $c$.
The quotient keeps precisely the non-$\ell^1$ parts and does not replace
such cancellation by an absolute tree majorant. This is why the theorem
can describe exceptional convergent parameters in a zero-gap system.
\end{remark}

\begin{corollary}[The top resonant parameters are unrestricted analytically]
A parameter of weight exactly $H$ does not occur in $\chi$. In particular,
if all resonant parameters have weight $H$, the candidates are either all
convergent or all divergent. This is a statement about the convergence
condition; the separate compatibility equations $\Psi=0$ still apply.
\end{corollary}

\begin{theorem}[Uniform convergence on compact parameter sets]
\label{thm:uniform}
For every compact subset $K\subset\{\chi=0\}$ there is $r_K>0$ such that
$Y(c)\in B_{r_K}$ for all $c\in K$, with uniformly bounded norms.
Locally on the algebraic locus $\{\chi=0\}$, this family is the restriction
of a holomorphic Banach-valued map. The assertion remains valid on singular
or reducible loci and, in particular, on $V(\Psi,\chi)$.
\end{theorem}
\begin{proof}
Choose actual coefficient families $\widetilde w_j\in\mathcal E$
representing the basis classes $w_j$ used above. Then
\[
 b_\nu=a_\nu-\sum_jq_{j\nu}\widetilde w_j\in\mathcal L.
\]
For $c$ satisfying $\chi(c)=0$, its window family equals
\begin{equation}\label{eq:window-polynomial}
 (Y_\gamma(c))_{\gamma\in J}=\sum_\nu b_\nu c^\nu.
\end{equation}
The right side is a finite polynomial with values in $\ell^1$.
It is uniformly bounded for $c\in K$. As a positive series its $B_r$ norm
tends to zero uniformly on $K$ as $r\downarrow0$.

The three inequalities \eqref{eq:tail-conditions} can therefore be imposed
with a common $r$. The resulting contraction supplies a common tail ball
and a common contraction constant. For the holomorphic assertion, use
\eqref{eq:window-polynomial} as a polynomial definition of $p(c)$ also
\emph{off} the locus. On a sufficiently small open parameter neighborhood
it obeys the same strict bounds. Iterating the tail map gives holomorphic
Banach-valued functions converging uniformly on smaller neighborhoods,
so its fixed point is holomorphic. On the locus, uniqueness identifies it
with the original formal candidate. Off the locus this extension need not
solve the original full equation, and no such claim is needed.
\end{proof}

\subsection{What is and is not effective}
The formal map $\Psi$ has finite input dependence. The map $\chi$ has
finite output dimension and bounded polynomial degree, but its coefficients
encode which finite linear combinations of the families $a_\nu$ lie in
$\ell^1$. These summability relations need not be decidable from an arbitrary
presentation of the input.

In particular, $\mathcal E/\mathcal L$ above is not being treated as a
Banach quotient, and no continuity of quotient-coordinate functionals is
assumed. Choosing a basis of its finite-dimensional algebraic subspace is
legitimate linear algebra. Replacing that choice by an effective exact
algorithm is a separate question with additional hypotheses.

\section{A sharp criterion for universal convergence}
\label{sec:gap}

Algebraicity describes a fixed equation even when some solutions diverge.
There is also a precise condition under which every admissible analytic
input has only convergent positive formal solutions. Here the relevant
numbers are eigenvalues of $A$, not the eigenvalue \emph{differences} in a
matrix-gauge problem.

Define
\begin{equation}\label{eq:gap}
 g(A,\G)=\inf\{\abs{\gamma-\lambda}:\gamma\in\G_{>0},\
 \lambda\in\spec(A),\ \gamma\ne\lambda\}.
\end{equation}

\begin{theorem}[Sharp universal convergence criterion]\label{thm:gap}
For fixed $A$ and $\G$, the following are equivalent:
\begin{enumerate}
\item $g(A,\G)>0$.
\item No positive real eigenvalue of $A$ is a left accumulation point of
$\G$.
\item For every $F$ satisfying \eqref{eq:F} and \eqref{eq:analytic-F},
every positive formal solution in $\HH_\G^m$ is convergent.
\end{enumerate}
If these conditions fail, an $F$ with a genuinely quadratic nonlinear
term and a divergent positive formal solution can be chosen arbitrarily
small in any prescribed joint norm of the form
\eqref{eq:analytic-F}.
\end{theorem}
\begin{proof}
Since $\G_{>0}$ has least element $\delta$, zero and negative real
eigenvalues cannot be approached. A nonreal eigenvalue has distance at
least the absolute value of its imaginary part from real exponents.
Only finitely many eigenvalues occur. Unbounded exponents cannot make
the distances in \eqref{eq:gap} tend to zero. Finally, no real number can
be a right accumulation point of a well-ordered real set: the part of
that set strictly above the number, if nonempty, has a least member.
It follows that failure of a positive gap is exactly left accumulation
at a positive real eigenvalue. This proves the first equivalence.

Assume $g(A,\G)>0$. The matrices $R_\gamma$ in
Section~\ref{sec:formal} then have a common finite norm bound $M$.
For nonresonant exponents this follows from boundedness of the matrix
resolvent on bounded closed sets separated from its poles and its decay
at infinity. At the finitely many exact resonances the chosen matrices
$R_\rho$ are fixed and finite. For any fixed $c$, the kernel polynomial
$C(x;c)$ tends to zero in $B_r$ as $r\downarrow0$.
The estimates \eqref{eq:smallness}--\eqref{eq:lipschitz} make
$y\mapsto C+\mathscr R F(x,y)$ a contraction on a sufficiently small
$B_r$ ball after shrinking $r$. Its fixed point equals $Y(c)$ by formal
uniqueness. Thus all candidates, and hence all compatible solutions,
converge.

For the converse let $\rho>0$ be a real eigenvalue approached from below.
Select strictly increasing $\alpha_n\in\G$ with
\begin{equation}\label{eq:bad-sequence}
 \rho/2<\alpha_n<\rho,
 \qquad 0<\rho-\alpha_n<2^{-2n}.
\end{equation}
Choose an eigenvector $v\ne0$ with $Av=\rho v$ and a linear functional
$\ell$ with $\ell(v)=1$; $\ell$ need not be a left eigenvector.
For positive constants $\varepsilon,\kappa$, put
\begin{equation}\label{eq:quadratic-bad}
 F(x,y)=\varepsilon\sum_{n\geq1}2^{-n}x^{\alpha_n}v
           +\kappa\ell(y)^2v.
\end{equation}
This satisfies the analytic majorant condition at every fixed positive
$x$-radius and finite $y$-radius. Taking $\varepsilon,\kappa$ small makes
its norm arbitrarily small while keeping its nonlinear term nonzero.

Seek a solution $y=u v$. The scalar equation is
\[
 (D-\rho)u=\varepsilon\sum_{n\geq1}2^{-n}x^{\alpha_n}+\kappa u^2.
\]
Set its optional free coefficient at $\rho$ to zero. Below $\rho$, all
nonzero forcing exponents exceed $\rho/2$, so the quadratic term
contributes nothing. At $\rho$ it contributes nothing either, and the
forcing has no term there. Hence a formal solution exists, including when
$\rho\notin\G$ so no compatibility equation is present. Its coefficients
are
\[
 u_{\alpha_n}=\frac{\varepsilon 2^{-n}}{\alpha_n-\rho},
 \qquad \abs{u_{\alpha_n}}>\varepsilon 2^n.
\]
Because the exponents stay bounded and tend to $\rho$, these coefficients
are not summable at any positive radius. The full vector $uv$ is an
actual formal solution of the matrix equation, regardless of the other
eigenvalues of $A$. This disproves the universal convergence assertion.
\end{proof}

\begin{remark}[Necessity is a universal, not pointwise, claim]
Zero gap does not imply that every equation, or every parameter choice,
diverges. The algebraic convergence locus may be empty, all of parameter
space, or a proper algebraic subset. The theorem constructs a bad input
when the gap fails; it does not replace the cancellation-sensitive
criterion for a particular input.
\end{remark}

\section{Universality: every affine algebraic set can occur}
\label{sec:universality}

The algebraicity theorem gives a substantial restriction, but it cannot
be strengthened to linearity or smoothness. The following construction
also separates formal solvability from convergence particularly cleanly.

\begin{theorem}[Universality of the analytic selection locus]
\label{thm:universality}
Let $P_1,\ldots,P_s\in\C[c_1,\ldots,c_n]$ be arbitrary polynomials.
There is a triangular system of the form \eqref{eq:system}, with convergent
Hahn coefficients and no formal compatibility obstruction, whose formal
solution parameter space is $\C^{n+s}$ and whose convergence locus is
\begin{equation}\label{eq:universal-locus}
 V(P_1,\ldots,P_s)\times\C^s.
\end{equation}
After normalizing the $s$ auxiliary resonant coefficients to zero, the
formal parameter space is $\C^n$ and the convergence locus is exactly
$V(P_1,\ldots,P_s)$.
\end{theorem}
\begin{proof}
Choose $h>0$ and an integer $d\geq\max(2,\deg P_1,\ldots,\deg P_s)$.
Set $\rho=(d+1)h$ and
\begin{equation}\label{eq:g-universal}
 g(x)=\sum_{r=2}^{\infty}\frac{x^{h-h/r}}{r^2},\qquad
 \G=\left\langle h,\ h-h/r\ (r\geq2)\right\rangle.
\end{equation}
The generating set is positive and well ordered, with least element $h/2$,
so $\G$ is admissible. The series $g$ is absolutely convergent at every
positive radius, because its exponents lie in $[h/2,h)$ and
$\sum r^{-2}<\infty$.

Write $P_j(c)=\sum_k p_{j,k}c^k$. Use variables $u_1,\ldots,u_n$ and
$v_1,\ldots,v_s$ and equations
\begin{align}
 Du_i&=h u_i,\qquad 1\leq i\leq n,\label{eq:univ-u}\\
 Dv_j&=\rho v_j+g(x)\sum_k p_{j,k}
               x^{(d-|k|)h}u^k,
           \qquad 1\leq j\leq s.\label{eq:univ-v}
\end{align}
All $x$-exponents in the nonlinear part are positive, and $d-|k|\geq0$.
Thus the structural conditions and the analytic majorant condition hold.
The residue matrix is diagonal with entries $h$ and $\rho$.

Every formal solution of \eqref{eq:univ-u} is $u_i=c_i x^h$.
Substitution in \eqref{eq:univ-v} gives
\[
 (D-\rho)v_j=P_j(c)\sum_{r=2}^{\infty}\frac{x^{\rho-h/r}}{r^2}.
\]
There is no forcing coefficient at the exact resonance $\rho$, so every
$c$ is formally allowed, and each $v_j$ has an independent free coefficient
$z_j$ at $\rho$. The full formal solutions are exactly
\begin{equation}\label{eq:univ-solution}
 u_i=c_i x^h,\qquad
 v_j=z_jx^\rho-\frac{P_j(c)}h
                  \sum_{r=2}^{\infty}\frac{x^{\rho-h/r}}r.
\end{equation}
At any positive radius, the terms $r_0^{\rho-h/r}$ tend to
$r_0^\rho>0$. The last series is therefore absolutely convergent exactly
when its scalar prefactor vanishes. This proves
\eqref{eq:universal-locus}. Normalizing $z_j=0$ removes the auxiliary
affine factor and changes no convergence condition.
\end{proof}

The ordinary-degree bound in Theorem~\ref{thm:algebraic} is also sharp.
Take one parameter of weight $h$ and a polynomial of exact degree $d$.
Here $H=(d+1)h$, so the bound is $d$, and the construction realizes a
degree-$d$ convergence equation. Choosing a polynomial with $d$ distinct
roots ensures that no lower-degree equation defines the same locus
in that parameter line.

\begin{example}[A two-point convergence locus]
Take $n=s=1$, $h=1$, $d=2$, and $P(c)=c^2-1$. The system is
\begin{equation}\label{eq:two-points}
 Du=u,\qquad Dv=3v+g(x)(u^2-x^2),\qquad
 g(x)=\sum_{r=2}^{\infty}r^{-2}x^{1-1/r}.
\end{equation}
All $(c,z)\in\C^2$ give formal solutions,
\[
 u=cx,\qquad v=zx^3-(c^2-1)\sum_{r=2}^{\infty}r^{-1}x^{3-1/r}.
\]
The convergent locus is exactly
$\{c=1\}\times\C\ \cup\ \{c=-1\}\times\C$.
Every coefficient of the equation is convergent; convergence of the
solution is selected by nonlinear cancellation, not by an input
summability failure.
\end{example}

\begin{example}[Singular and reducible convergence loci]
Using $P(c_1,c_2)=c_1c_2$ gives the union of the two coordinate axes.
Using $P(c_1,c_2)=c_2^2-c_1^3$ gives a cusp. These examples show that the
word ``algebraic'' in Theorem~\ref{thm:algebraic} does not hide an
unproved regularity assertion. The theorem concerns the complex point
set; no claim about a canonical nonreduced scheme structure is intended.
\end{example}

\subsection{Formal compatibility can also encode arbitrary varieties}
For comparison, omit the accumulating multiplier $g$ and use
\begin{equation}\label{eq:formal-univ}
 Du_i=h u_i,\qquad
 Dv_j=dh v_j+\sum_kp_{j,k}x^{(d-|k|)h}u^k.
\end{equation}
Now $u=cx^h$ gives exact-resonance forcing $x^{dh}P_j(c)$.
The formal compatibility equations are $P_j(c)=0$, and the formal solution
space is $V(P)\times\C^s$. All those solutions converge, by local
finiteness. Thus the same algebraic set can arise either as a formal
resonance obstruction or as a purely analytic selection condition. The
accumulating construction proves that the two mechanisms are not the same.

\section{Nonlinear generation of a sharp small-divisor threshold}
\label{sec:hidden}

The universal gap criterion uses the whole exponent monoid. It cannot in
general be replaced by separation of the original forcing support. The
following scalar example gives exact formulas and a complete threshold,
including compatibility at the resonance.

Let
\begin{equation}\label{eq:hidden-support}
 \alpha_n=2-\frac1n\quad(n\geq3),\qquad
 \G=\langle1,\alpha_3,\alpha_4,\ldots\rangle.
\end{equation}
The least positive exponent is $1$, and the only members of $\G$ up to
$3$ are
\begin{equation}\label{eq:window-support}
 0,1,2,3,\qquad \alpha_n,\qquad 1+\alpha_n\quad(n\geq3).
\end{equation}
Indeed, two $\alpha$-generators have sum at least $10/3>3$.
Below $3$, an integer plus an $\alpha$-generator can use at most one $1$.
In particular,
\begin{equation}\label{eq:hidden-ancestry}
 \Anc(3)=\{0,1,2,3\},
\end{equation}
although infinitely many exponents lie below $3$.

For $p>1$ and $a,b\in\C$ consider
\begin{equation}\label{eq:hidden-system}
 (D-3)y=a x+\sum_{n=3}^{\infty}n^{-p}x^{\alpha_n}
                  +b y^2+2b^2y^3.
\end{equation}
Its forcing is absolutely convergent at every positive radius.
Each original nonzero forcing exponent is at distance strictly greater
than $1$ from the spectral value $3$; their infimum distance is $1$.
There is no original forcing at the resonance.

\begin{theorem}[Exact nonlinear threshold]\label{thm:hidden}
Every $a,b$ and every free coefficient $c=y_3$ produce a unique positive
formal solution of \eqref{eq:hidden-system}. Its coefficients in the
critical window are
\begin{align}
 y_1&=-\frac a2,&y_2&=-\frac{a^2b}{4},\label{eq:hidden-coef1}\\
 y_{\alpha_n}&=-\frac{n^{1-p}}{n+1},&
 y_{1+\alpha_n}&=-\frac{ab\,n^{2-p}}{n+1},\label{eq:hidden-coef2}\\
 y_3&=c.\label{eq:hidden-coef3}
\end{align}
The formal solution is convergent if and only if
\begin{equation}\label{eq:hidden-threshold}
 p>2\quad\text{or}\quad ab=0.
\end{equation}
The criterion is independent of the top resonant parameter $c$.
\end{theorem}
\begin{proof}
At exponents $1$ and $2$, coefficient comparison gives
$y_1=-a/2$ and $y_2=-b y_1^2=-a^2b/4$.
The exponent $\alpha_n$ is below $2$, so no nonlinear term contributes:
\[
 y_{\alpha_n}=\frac{n^{-p}}{\alpha_n-3}
             =-\frac{n^{1-p}}{n+1}.
\]
At $1+\alpha_n=3-1/n$, the only nonlinear contribution is
$2b y_1y_{\alpha_n}$. The cubic term starts at exponent $3$.
Dividing by $-1/n$ gives the second coefficient in
\eqref{eq:hidden-coef2}.

At the exact resonance $3$, only the ancestors in
\eqref{eq:hidden-ancestry} can enter. The forcing is
\[
 2b y_1y_2+2b^2y_1^3
 =\frac{a^3b^2}{4}-\frac{a^3b^2}{4}=0.
\]
Hence compatibility holds for every $a,b$ and $c=y_3$ is free.
Theorem~\ref{thm:formal} supplies the unique formal continuation.

The first infinite coefficient family in \eqref{eq:hidden-coef2} is
summable for every $p>1$. The second has magnitude comparable to
$|ab|n^{1-p}$. It is summable exactly when $ab=0$ or $p>2$.
All exponents in this calculation are bounded by $3$, so no radius can
alter this summability test. Theorem~\ref{thm:window}, with $H=3$,
proves that this is also sufficient for the full solution to converge.
\end{proof}

\begin{remark}[The two roles of cancellation]
The cubic term is chosen to cancel the \emph{exact} resonance, making all
formal branches exist. It does not cancel the infinitely many quadratic
near-resonances below $3$. Thus formal compatibility alone cannot detect
the threshold. Conversely, setting $a=0$ or $b=0$ removes the entire
non-summable family, illustrating why an absolute path bound can miss a
convergent specialization.
\end{remark}

\subsection{Finite coefficient data cannot determine the analytic locus}

\begin{proposition}[No general finite-coefficient convergence certificate]
\label{prop:no-finite}
For a suitable fixed $A$ and $\G$, and any prescribed finite set of input
coefficient positions, there exist two analytic inputs agreeing at all
those positions and having identical formal compatibility maps, while one
has convergent positive formal solutions and the other has none.
The inputs can be arbitrarily close in the majorant norm.
\end{proposition}
\begin{proof}
Take the scalar $A=1$ and
$\G=\langle1,1-1/n\ (n\geq2)\rangle$.
Use $F^{(0)}=0$ and, after discarding a sufficiently long finite prefix,
\[
 F^{(1)}(x,y)=\varepsilon\sum_{n>N}n^{-2}x^{1-1/n}.
\]
Choose $N$ so that none of the displayed coefficient positions is among
the finitely prescribed ones. Both inputs have zero forcing at the exact
resonance $1$, hence their compatibility maps are identically zero.
The first has solutions $cx$. The second has solutions
\[
 cx-\varepsilon\sum_{n>N}n^{-1}x^{1-1/n},
\]
none of which is absolutely convergent. The norm of the difference tends
to zero with $\varepsilon$, or with $N\to\infty$ at fixed radius.
\end{proof}

This does not contradict the finite polynomial description of the analytic
locus. In the example the extra polynomial equation is simply a nonzero
constant, so the locus is empty. Determining that constant obstruction
requires information about an infinite coefficient tail.

\section{Computation, proof audit, and a formalization route}
\label{sec:verification}

\subsection{What the supplied program checks}
The program \path{verification/verify.py} uses exact rational exponents,
exact SymPy matrices, and symbolic polynomial coefficients. Its central
routine enumerates a \emph{finitely generated rational} monoid up to a
cutoff and implements \eqref{eq:recursion}. At a singular $T_\rho$, it
constructs a right inverse by pivot columns and rows, rather than assuming
that kernel and cokernel projections coincide.

After constructing each candidate, the program independently reevaluates
$F$ and checks the residual identity \eqref{eq:residual} at every enumerated
exponent. Thus an incompatible candidate is expected to have exactly the
recorded resonant residual, not a spuriously vanishing one.

The recorded run used Python 3.13.5 and SymPy 1.14.0, with seed $20260929$.
It produced the following actual counts:
\begin{center}
\begin{tabular}{lr}
\toprule
Check category & Count\\
\midrule
Finite nonlinear systems & 21\\
Exact scalar equalities, including matrix entries & 670\\
Sharp Riccati coefficient identities in the independent recurrence & 24\\
Seeded additional matrix systems & 12\\
\bottomrule
\end{tabular}
\end{center}
The last two rows are included in, not added to, the equality and system
counts. The systems include six sharp-degree examples, the generated
small-divisor example, the arbitrary-variety construction, a Jordan-block
example, and twelve seeded systems.

For the sharpness family, the first coefficients of the obstruction after
removing $a^n b^{n-1}\delta^{-2n+2}$ are
\begin{center}
\begin{tabular}{ccl}
\toprule
$n$ & Coefficient & Total degree\\
\midrule
2 & $1$ & 3\\
3 & $1/4$ & 5\\
4 & $1/36$ & 7\\
5 & $1/576$ & 9\\
6 & $1/14400$ & 11\\
7 & $1/518400$ & 13\\
\bottomrule
\end{tabular}
\end{center}
These agree exactly with \eqref{eq:sharp-obstruction}.

\subsection{A Jordan-block check}
Take
\[
 A=\begin{pmatrix}1&1\\0&1\end{pmatrix},\qquad
 F=x\begin{pmatrix}b\\a\end{pmatrix}.
\]
At exponent $1$, $T_1$ has image spanned by the first coordinate and
kernel also spanned by the first coordinate, but the cokernel complement
is the second coordinate. The correct compatibility condition is $a=0$,
and the candidate coefficient is $(c,-b)^{\mathsf T}$. Projecting onto
the kernel in place of the cokernel would give the wrong condition.
This exact test exercises a distinction on which the general theorem
depends.

\subsection{Limits of the computation}
The code does not prove that an arbitrary infinite set is well ordered,
that an infinite coefficient family is summable, or that the algebraic
quotient construction is effective. Restricting an infinite example to
finitely many generators verifies its displayed coefficient identities,
not its convergence threshold. The threshold and universality arguments
are mathematical proofs in Sections~\ref{sec:universality} and
\ref{sec:hidden}.

No Lean theorem is supplied or claimed for the new results. The repository's
\path{NeumannWords.lean} was inspected as an available support-theory
component, not recompiled in this work. The numerical and symbolic checks
must not be described as a proof-assistant verification.

\subsection{A useful implementation contract}
For effective finite data, the core formal algorithm is:
\begin{enumerate}
\item Determine the exact real resonances and certify their finite ancestor
sets. Choose kernel and cokernel splittings over the coefficient field.
\item Process only those ancestor exponents needed by the requested
certificate. Compute $f_\gamma$, then $R_\gamma f_\gamma+C_\gamma$, and
record $Q_\rho f_\rho$ at a resonance.
\item For convergence, separately represent the window coefficient
families and certify their linear relations modulo $\ell^1$. Only after
that step construct the analytic equations $\chi$.
\end{enumerate}
The third step cannot generally be replaced by a finite coefficient cutoff,
as Proposition~\ref{prop:no-finite} shows. An infinite-support API must
expose exact ancestry and summability certificates, rather than pretending
that ordinary bounded enumeration always terminates.

\subsection{Formalization in layers}
A Lean development can keep the difficult parts separate. The support layer
can reuse positive-word and antidiagonal lemmas already present in ProveIt.
The finite-dimensional layer needs right inverses, kernel/cokernel
splittings, and the coefficient residual identity. The formal layer proves
the uniform valuation gain and finite polynomial ancestry.

The analytic layer then builds the weighted $\ell^1$ algebra, coefficient
projections, resolvent bounds, and the tail contraction theorem. The central
algebraicity step is a finite-dimensional linear-algebra argument in the
quotient of all window families by the summable subspace. It requires no
computability claim and no topology on that quotient. The uniform-parameter
theorem additionally needs holomorphic dependence of a uniformly contractive
Banach-space fixed point.

A sensible first formal target is Theorem~\ref{thm:window}; it isolates the
new analytic reduction and is independent of the choice of quotient basis.
The universality construction is a second tractable target because all its
coefficients and summability decisions are explicit.

\section{Further research questions and concrete next targets}
\label{sec:future}

\subsection{Effective summability quotients}
Theorem~\ref{thm:algebraic} reduces analytic classification to linear
relations among finitely many coefficient families modulo $\ell^1$.
For which exact sequence classes is this relation problem decidable?
A first target is a family whose coefficients admit certified expansions
in $n^\alpha(\log n)^\beta$ with summable remainders. One should prove
that cancellation of the finitely many non-summable asymptotic types is
both necessary and sufficient, then extract $\chi$ algorithmically.
Oscillatory families require absolute, not merely conditional, summability.

\subsection{Minimal analytic-obstruction dimension}
The proof bounds the number of equations by a monomial count, which can be
very loose. Determine the possible dimensions of the quotient span $W$ for
a prescribed spectrum and coupling pattern. A concrete target is to
classify two-component quadratic systems with two positive resonances,
including when one obstruction polynomial suffices. The relevant invariant
should retain cancellations rather than count all substitution trees.

\subsection{Polynomial families of equations}
If finitely many external parameters enter the coefficients with uniformly
bounded polynomial degree, a bounded exponent window again has a finite
polynomial template. Formulate a joint parameter theorem including both
input parameters and resonant solution parameters, under a common analytic
majorant on compact sets. Sharp degree bounds should account separately
for positive exponent cost and zero-cost parameter dependence. Without a
polynomial hypothesis, an algebraic conclusion is not expected.

\subsection{Moving eigenvalues and wall crossing}
Our matrix $A$ is fixed. When its eigenvalues vary, the resonant parameter
space changes and the coefficient inverses become rational functions of
the spectral parameters. Can one obtain algebraic or constructible
convergence loci on a suitable stratification away from moving accumulation
walls? The scalar multiplier $(\gamma-\lambda)^{-1}$ already shows why
one cannot simply repeat the finite-degree proof with $\lambda$ treated as
another polynomial parameter.

\subsection{Logarithmic solutions after compatibility failure}
The current theory deliberately excludes logarithms. If $\Psi(c)\ne0$,
which nonlinear systems can be solved after adjoining a bounded degree in
$\log x$, and when does the nonlinearity force an unbounded logarithmic
expansion? A first test family is the sharp Riccati example
\eqref{eq:sharp}, whose linearization gives explicit resonant behavior.
The goal is a theorem distinguishing a finite log correction from a
change of asymptotic scale.

\subsection{Higher-rank exponent groups}
The finite-degree argument uses the real bound $N\delta\leq H$ to bound
$N$. It need not survive in a non-Archimedean ordered group, even though
classical Neumann word finiteness remains available for each exact
exponent. Find a substitute for a bounded real cost window that still
bounds the number of parameter monomials relevant to analytic failure.
A useful first result would be either a rank-two positive theorem under a
cofinality hypothesis or an explicit nonalgebraic parameter locus.

\subsection{Irregular equations and unbounded dangerous exponents}
For $x^{k+1}y'=Ay+F(x,y)$ with $k>0$, coefficient transport is no longer
purely diagonal in the exponent. Classical factorial divergence can
persist at arbitrarily large exponents, so the bounded-window argument
does not apply directly. Determine whether the parameter locus of a fixed
Gevrey or Borel-summability class is still algebraic in a carefully chosen
irregular model. This is a separate problem, not a consequence of the
present theorem.

\subsection{Summation beyond absolute convergence}
The divergent candidates in the universality construction may still have
renormalized analytic realizations. Classify realizations compatible with
multiplication and $D$, and identify which normalization data are invisible
to finite asymptotic prefixes. In that enlarged class the selected parameter
locus could differ from $\cV_{\mathrm{an}}$. Absolute convergence must
remain explicit in every comparison.

\subsection{Geometry, stability, and certified computation}
The convergence locus is algebraic for a fixed equation but can change
under arbitrarily small input perturbations. Characterize subclasses in
which its dimension or defining ideal is stable. A practical first target
is a finite collection of certified non-summable asymptotic types with
uniformly summable remainders; the problem may then reduce to finite matrix
rank stability. Such a theorem would give meaningful perturbation
certificates rather than an unjustified generic robustness claim.

\subsection{A complete machine-checked model case}
Formalize the two-point system \eqref{eq:two-points}, including positivity
and well ordering of its generated support, exact formal coefficients,
and the equivalence of convergence with $c^2=1$. This would provide a
compact end-to-end test of support algebra, coefficient recursion, absolute
summability, and parameter cancellation. The proof should be independent
of numerical truncation.

\section{Conclusion}

The nonlinear problem has two separate algebraic geometries. Exact
resonances give a finite-input compatibility variety. Absolute convergence
gives another algebraic condition, obtained from possibly infinite
coefficient families rather than from a finite input jet.

The decisive analytic fact is that, for positive real Hahn exponents in a
regular-singular Euler equation, every possible divergence lies below the
largest positive real eigenvalue. Above that cutoff the resolvent is
bounded and the nonlinear tail is contractive after shrinking the radius.
Positive parameter weights make the remaining infinite window a finite
polynomial family. Passing to the algebraic quotient by $\ell^1$ therefore
produces finitely many convergence equations.

The universality construction shows both the strength and the limit of
this classification: the convergence locus cannot be an arbitrary
nonalgebraic subset, but it can be an arbitrary affine algebraic set.
The generated small-divisor example further shows that original forcing
separation and exact-resonance cancellation do not by themselves imply
convergence. These results provide a proof-based extension of the
repository's linear accumulation analysis, with the unresolved effective,
logarithmic, and higher-rank questions stated separately.

\appendix
\section{Proof dependencies and hypothesis audit}
\label{app:audit}

\begin{longtable}{@{}p{0.29\textwidth}p{0.65\textwidth}@{}}
\toprule
Result & Essential hypotheses and dependency\\
\midrule
\endhead
Formal parametrization & Positive real monoid; well ordering; no constant
or zero-exponent linear term in $F$. Analytic convergence is not used.\\[4pt]
Finite ancestry & Exact additive divisors, not finiteness of a bounded
interval. The real lower bound supplies an explicit degree bound.\\[4pt]
Tail convergence & Joint analytic majorant; fixed finite matrix; positivity
of the unknown; well-order right gap above the final real eigenvalue. No
left spectral gap is needed.\\[4pt]
Algebraic analytic locus & Finite resonant parameter set with positive
weights; bounded-window criterion; algebraic quotient by the $\ell^1$
subspace. No topological quotient or effective basis is assumed.\\[4pt]
Uniform radii & A finite polynomial representation by summable coefficient
families on the locus, followed by a common contraction estimate.\\[4pt]
Universal gap converse & Quantification over every admissible $F$; an
invariant eigenvector direction and a selected left-accumulating sequence.
The constructed nonlinear formal solution exists.\\[4pt]
Universality & A fixed well-ordered accumulating monoid and polynomial
coupling in the unknowns; all input coefficients are absolutely summable.
Auxiliary free resonant coefficients account for the affine factor.\\
\bottomrule
\end{longtable}

\subsection{Frequently tempting but invalid shortcuts}
A bounded set of Hahn exponents need not be finite. A formal coefficient
recursion does not imply absolute convergence. An eigenvector is not
necessarily paired with a left eigenvector having nonzero pairing, so the
universal counterexample uses only an arbitrary functional $\ell(v)=1$
and explicitly constructs a solution in the invariant line. A kernel
projection does not replace a cokernel obstruction for a Jordan block.
The quotient by summable families is used algebraically, without an
unjustified continuity assertion. Finally, finite algebraicity of a
parameter locus does not mean finite-coefficient decidability of that locus.

\section{Repository provenance and literature boundary}
\label{app:provenance}

The source comparison used the following pinned commit:
\begin{center}\snap\end{center}
In particular, the linear Hahn--Fuchsian package was read in its README, introductory
framework, explicit accumulation model, verification discussion, and
further-research sections. The support module
\path{Algebra/SurrealNumbers/Surreal/HahnSeries/NeumannWords.lean}
was inspected through its definitions and principal word-finiteness proof.
Repository searches also checked related resonance and convergence-locus
terminology. This is not a line-by-line audit of the entire repository or
its large canonical transseries volume.

The cited Gontsov--Goryuchkina sources were checked for their stated
convergence setting and finite-generation framework. Their results,
van der Hoeven's generalized-series fixed-point theory, and existing Hahn
convergence definitions are acknowledged as prior theory, not relabeled as
new contributions. The proof of the central algebraic-locus theorem is
self-contained, but independence of historical discovery remains a matter
for further specialist review.

The delivered package includes the standalone source, compiled PDF,
reproduction script, recorded exact results, a source audit, and build
instructions. No repository files were modified. No private data or
external font files are required to rebuild the article.

\clearpage
\begin{thebibliography}{9}
\raggedright
\bibitem{repo-linear}
VladimirReshetnikov/ProveIt,
\emph{Finite Resonance Certificates and Sharp Analytic Normalization for
Hahn--Fuchsian Systems}, research package, 29 September 2026, with repository
editorial amendments. Inspected at commit \snap.
\url{https://github.com/VladimirReshetnikov/ProveIt/tree/3d5973524506411392a911470b5ddc35521568ea/Analysis/Transseries/docs/series-and-transseries/Hahn_Fuchsian_Resonance_Analytic_Normalization}.

\bibitem{repo-words}
VladimirReshetnikov/ProveIt,
\emph{NeumannWords.lean}, module on positive ordered-word support finiteness
and strongly summable labeled products. Inspected at commit \snap.
\url{https://github.com/VladimirReshetnikov/ProveIt/blob/3d5973524506411392a911470b5ddc35521568ea/Algebra/SurrealNumbers/Surreal/HahnSeries/NeumannWords.lean}.

\bibitem{vdh}
J. van der Hoeven,
\emph{Operators on generalized power series},
Illinois Journal of Mathematics \textbf{45} (2001), 1161--1190.
\url{https://doi.org/10.1215/ijm/1258138061}.
Author version: \url{https://www.texmacs.org/joris/noeth/noeth.html}.

\bibitem{gg2015}
R. R. Gontsov and I. V. Goryuchkina,
\emph{On the convergence of generalized power series satisfying an
algebraic ODE}, Asymptotic Analysis \textbf{93} (2015), 311--325.
\url{https://doi.org/10.3233/ASY-151297}.
Preprint: \url{https://arxiv.org/abs/1407.1330}.

\bibitem{gg2024}
R. Gontsov and I. Goryuchkina,
\emph{Convergence of (generalized) power series solutions of functional
equations}, arXiv:2412.00778 (2024), extended abstract and manuscript.
\url{https://arxiv.org/abs/2412.00778}.
English extended abstract: \url{https://arxiv.org/html/2412.00778v1}.

\bibitem{js2023}
M. Joswig and B. Smith,
\emph{Convergent Hahn series and tropical geometry of higher rank},
Journal of the London Mathematical Society \textbf{107} (2023), 1450--1481.
\url{https://doi.org/10.1112/jlms.12716}.
\end{thebibliography}
\end{document}
